\documentclass{osucourse}
\course{4581}

\begin{document}

\CatalogDescription
\PrerequisitesAndExclusions

\begin{Purpose}
Math 4580--4581 constitutes a two-semester sequence on abstract algebra,
intended to introduce students the main concepts of this subject area.
Focused on groups, rings and fields, this course gives the students a
deep understanding of these three basic algebraic structures, and
provides a good foundation for more specialized work. A significant goal
of the course is to improve mathematical reasoning and proof writing.

Math 4581 presents special classes of groups, group actions on sets,
vector spaces, and field theory that concludes with elements of Galois
theory. The course places these topics in their historical context where
possible.
\end{Purpose}

\begin{Text}
Tom Judson, \emph{Abstract Algebra: Theory and Applications}

\url{https://aimath.org/textbooks/approved-textbooks/judson/}
\end{Text}

\begin{TopicsOutline}[ (Chapters from Judson book)]
12. \textbf{Matrix Groups and Symmetry}

12.1. Matrix Groups

12.2 Symmetry

13. \textbf{The Structure of Groups}

13.1. The Structure of Finite Abelian Groups

13.2. Solvable Groups

14. \textbf{Group Actions on Sets}

14.1. Group Actions

14.2. The Class Equation

14.3. Burnside's Theorem

20. \textbf{Vector Spaces} (over arbitrary fields)

20.1. Definitions and Examples

20.2. Subspaces

20.3. Independence

17. \textbf{Polynomials} (a quick 1 day review of the results we will
need for field theory)

21. \textbf{Fields}

21.1. Extensions

21.2. Splitting Fields

21.3. Geometric Constructions

22. \textbf{Finite Fields} (this chapter contains both characteristic
and separability; cover them first)

22.1. Finite fields

23. \textbf{Galois Theory}

23.1. Field Automorphisms

23.2. Fundamental Theorem of Galois Theory

23.3. Applications (insolvability of the quintic by radicals)
\end{TopicsOutline}

\begin{SuggestedSchedule}
\week{1}{12.1--12.2 Matrix groups and symmetry.}
\week{2}{13.1 Structure of finite abelian groups.}
\week{3}{13.2 Solvable groups; composition series.}
\week{4}{14.1 Group actions on sets.}
\week{5}{14.2 The class equation; the Sylow theorems.}
\week{6}{14.3 Burnside's Theorem and counting applications. \emph{Midterm I.}}
\week{7}{20.1--20.2 Vector spaces over arbitrary fields; subspaces.}
\week{8}{20.3 Independence, bases, and dimension.}
\week{9}{17 Review of polynomials; 21.1 field extensions.}
\week{10}{21.1 Algebraic and finite extensions; degree.}
\week{11}{21.2--21.3 Splitting fields; geometric constructions. \emph{Midterm II.}}
\week{12}{22.1 Finite fields; characteristic and separability.}
\week{13}{23.1 Field automorphisms.}
\week{14}{23.2--23.3 Fundamental Theorem of Galois Theory; insolvability of the quintic.}
\finalsweek{Final examination.}
\end{SuggestedSchedule}

\end{document}
